\paragraph{Buffer manager}
Databases use locks (used to avoid conflicting updates to the data) and latches (avoid conflicting updates to data structures of system) to manage buffers.

The buffer cache latches are used to avoid conflicting accesses to hash buckets with the block headers and cover several hash buckets, the number of which is configurable.
The reason that there isn't a latch per bucket, or even per block header is to reduce the space devoted to the engine data structures.

Of course, as with any other kind of locking, only one process can hold a latch and they typically span several linked lists of block headers
and the typical performance bottlenecks also apply.

To mitigate these performance issues, we can
\begin{itemize}
    \item reduce the number of data in a block
    \item Configure a DB with more latches and less buckets per latch
    \item Multiple buffer pools
    \item Tune queries to reduce full table scans
    \item Avoid many concurrent queries that access the same data
    \item Avoid concurrent transactions
\end{itemize}

The link list in which a block header resides is found using hashing on some kind of block header.
Then, the linked list is traversed to find the block, which is expensive, which is why the length of these lists should be reduced as much as possible.

Block headers contain the following data:
\begin{itemize}
    \item Block number or ID
    \item Format
    \item LSN (Log Sequence Number), can be similar to git commit refs, sometimes also simply a counter
    \item Checksum for integrity
    \item Latches / status / etc flags
    \item Buffer replacement information (for replacement policies)
\end{itemize}

Status information of the block includes relevant data for management of the buffer, such as pinning (i.e. pinning it into memory), usage counts, clean / dirty flag
and more.
